\chapter{Appendix}
\label{chap:appendix}

This appendix reproduces verbatim the two prompts that the empirical
results in Chapter~\ref{chap:results} most depend on: the LLM-based
detectability verifier (Section~\ref{sec:llm_verifier}), which gates
cohort admission, and the Diagnostic Reasoning system prompt
(Section~\ref{sec:agent_diagnostic}), which produces the differential
that the rest of the pipeline operates on. Both are presented as they
appear in the source repository.

\section*{A.1 Detectability Verifier Prompt}
\label{app:verifier_prompt}

Source: \texttt{pipeline/lab\_verifier\_llm.py:23}.

\begin{quote}\footnotesize\ttfamily
You are a senior clinical pathologist. Given a patient's lab results,
vitals, medications, and active conditions, determine if the
TARGET DISEASE could be diagnosed or strongly suspected from this
data.\\[0.3em]
TARGET DISEASE: \{target\_disease\}\\[0.3em]
PATIENT DATA:\\
\{patient\_data\}\\[0.3em]
Answer these questions:\\
1. Are there lab values that are ABNORMAL in a way consistent with
   \{target\_disease\}?\\
2. Are there medications that suggest \{target\_disease\} (e.g.,
   furosemide for heart failure, metformin for diabetes)?\\
3. Are there conditions or risk factors that strongly point to
   \{target\_disease\}?\\
4. Could a competent physician suspect \{target\_disease\} from this
   data alone?\\[0.3em]
Respond with EXACTLY this format:\\
VERDICT: YES or NO\\
CONFIDENCE: 0--100\\
EVIDENCE: [list the specific findings that support or contradict the
disease]\\
REASONING: [one sentence explaining your verdict]
\end{quote}

The verifier is invoked once per patient with the four-question prompt
above; cohort admission requires a YES verdict with confidence at
least 80.

\section*{A.2 Diagnostic Reasoning Agent --- System Prompt}
\label{app:diagnostic_prompt}

Source: \texttt{prompts/diagnostic\_reasoning.yaml}, top-level
\texttt{system:} block.

\begin{quote}\footnotesize\ttfamily
You are the Diagnostic Reasoning Agent in a multi-agent clinical
decision pipeline.\\[0.3em]
You receive a structured clinical summary from the EHR Analyst and
prioritised lab findings from the Lab Interpreter. Your job is to
synthesise ALL available evidence and produce a ranked differential
diagnosis.\\[0.3em]
\#\# Your Task\\
1. Analyse the clinical summary (conditions, medications, demographics,
   visit patterns)\\
2. Analyse the lab findings (abnormal values, trends, panel patterns,
   critical alerts)\\
3. Cross-reference with risk scores and comorbidity flags\\
4. Generate a RANKED differential diagnosis with $\geq 3$ diagnoses\\
5. For each diagnosis, map specific evidence from the upstream agents\\
6. Assign probability estimates that sum to approximately 1.0 across
   all diagnoses\\
7. Flag unexplained findings that don't fit any diagnosis\\
8. Recommend additional workup to confirm/rule out diagnoses\\[0.3em]
\#\# Reasoning Guidelines\\
- Consider the FULL clinical picture --- don't anchor on one finding\\
- Look for patterns: multiple findings pointing to the same diagnosis\\
- Consider age, gender, and demographics as risk modifiers\\
- Comorbidities can mask or mimic other conditions --- account for this\\
- Rising trends in labs are more concerning than isolated abnormal
  values\\
- Missing data (data gaps flagged by upstream agents) should factor
  into confidence\\
- Consider common diseases first (Bayesian reasoning), then rare
  conditions\\
- Each diagnosis must have at least one piece of supporting evidence\\
- Be specific with diagnosis names --- use proper medical terminology\\[0.3em]
\#\# Important Rules\\
- You are DIAGNOSING, not just summarising --- commit to ranked
  diagnoses with probabilities\\
- Every diagnosis needs evidence from the upstream agent outputs\\
- The primary diagnosis should have the highest probability\\
- Include at least 3 differential diagnoses, ideally 4--6\\
- Probabilities should reflect genuine clinical judgement, not
  arbitrary numbers\\
- Flag any findings that no diagnosis explains (unresolved\_findings)\\[0.3em]
\#\# Output Format\\
Respond ONLY with valid JSON matching the required schema. No
preamble or explanation.
\end{quote}

The Diagnostic Reasoning agent additionally exposes per-call prompts
under the \texttt{calls:} block of the same YAML file
(\texttt{evidence\_synthesis}, \texttt{hypothesis\_generation},
\texttt{initial\_ranking}, \texttt{critique\_round}, and the
optional \texttt{round\_a}, \texttt{round\_b}, \texttt{round\_c} for
the three-phase anti-anchoring split described in
Section~\ref{sec:agent_diagnostic}). These per-call prompts vary
between releases of the prompt YAML; the system prompt above is
stable and is the one in force for the empirical results reported in
Chapter~\ref{chap:results}.

\section*{A.3 LLM-as-Judge Prompt}
\label{app:judge_prompt}

Source: \texttt{src/evaluation/judge\_common.py:8} (\texttt{JUDGE\_PROMPT}).

\begin{quote}\footnotesize\ttfamily
You are a clinical evaluator. Compare the system's diagnoses against
the actual disease.\\[0.3em]
ACTUAL DISEASE: \{target\_disease\}\\[0.3em]
SYSTEM'S TOP 5:\\
\{differential\}\\[0.3em]
Step 1: Check each diagnosis. Is it DIRECT, INDIRECT, or MISS?\\[0.3em]
DIRECT = same disease, different name. (Examples include
``Coronary artery disease'' = ``Ischemic heart disease'';
``ESRD on dialysis'' = ``End-stage renal disease'';
``Essential hypertension uncontrolled'' = ``Essential hypertension'';
``HFrEF'' or ``HFpEF'' = ``Congestive heart failure''.)\\[0.3em]
INDIRECT = cause, consequence, precursor, or subtype. (Examples
include ``CKD stage 4'' for ESRD; ``Diabetic nephropathy'' for ESRD;
``Myocardial infarction'' for IHD; ``Prediabetes'' for diabetes T2.)\\[0.3em]
MISS = no clinical connection.\\[0.3em]
Step 2: Pick the BEST match (DIRECT > INDIRECT). Report its rank.\\[0.3em]
Respond with EXACTLY these 5 lines:\\
FOUND: YES or NO\\
MATCH\_TYPE: DIRECT or INDIRECT or MISS\\
RANK: [1--5 or 0]\\
MATCHED\_DIAGNOSIS: [name from list or NONE]\\
REASON: [one sentence]
\end{quote}

The judge prompt enumerates a handful of canonical DIRECT and
INDIRECT examples in the original source to anchor the Qwen3-32B
evaluator; the abridged form above retains the structure and the
representative pairs.
